\documentclass[11pt,a4paper]{article}
\usepackage[utf8]{inputenc}
\usepackage[T1]{fontenc}
\usepackage[french]{babel}
\usepackage{amsmath,amssymb}
\usepackage[top=2cm,bottom=2.5cm,left=2.5cm,right=2cm]{geometry}
\usepackage{enumitem}
\usepackage{booktabs}

\begin{document}

\begin{center}
{\large\bfseries Préparation au concours d'entrée à l'IFORD -- Yaoundé}\\[2mm]
{\Large\bfseries Concours TYPE B -- Épreuve de PROBABILITÉS ET STATISTIQUE}\\[2mm]
Sujet d'entraînement n\textsuperscript{o}\,1 \quad --\quad Durée : 4 heures\\[1mm]
\emph{Sujet rédigé pour l'entraînement, conforme au programme officiel. Ce n'est pas un sujet officiel.}
\end{center}
\hrule\medskip

\noindent\textbf{Consignes.} Les exercices sont indépendants. Calculatrice non programmable autorisée.
On donne : $u_{0{,}95} = 1{,}645$ ; $u_{0{,}975} = 1{,}96$ ; $e^{-2} \approx 0{,}1353$ ; $2^{0{,}4} \approx 1{,}3195$.

\section*{Exercice 1 -- Distribution à deux variables (5 points)}
Population urbaine (en millions) d'un pays lors de cinq recensements successifs, numérotés $x$ :
\begin{center}
\begin{tabular}{lccccc}
\toprule
$x_i$ & 1 & 2 & 3 & 4 & 5\\ \midrule
$y_i$ & 2{,}1 & 2{,}9 & 4{,}2 & 4{,}8 & 6{,}0\\ \bottomrule
\end{tabular}
\end{center}
\begin{enumerate}
  \item Calculer $\bar{x}$, $\bar{y}$, $\operatorname{Var}(x)$, $\operatorname{Var}(y)$ et $\operatorname{Cov}(x,y)$.
  \item Déterminer la droite de régression de $y$ en $x$ par la méthode des moindres carrés.
  \item Calculer le coefficient de corrélation linéaire $r$ et l'interpréter.
  \item Prévoir la population urbaine au $7^{\text{e}}$ recensement. Quelle réserve faut-il émettre ?
\end{enumerate}

\section*{Exercice 2 -- Estimation et test (5 points)}
Dans une enquête démographique et de santé, sur un échantillon aléatoire de $400$ femmes,
$92$ déclarent utiliser une méthode contraceptive moderne.
\begin{enumerate}
  \item Donner une estimation ponctuelle de la prévalence contraceptive $p$.
  \item Construire un intervalle de confiance de $p$ au niveau $95\,\%$.
  \item Les autorités affirment que la prévalence dépasse désormais $20\,\%$. Tester au seuil de $5\,\%$
  $H_0 : p = 0{,}20$ contre $H_1 : p > 0{,}20$. Conclure.
  \item Quelle taille d'échantillon faudrait-il pour que l'intervalle de confiance à $95\,\%$ ait une demi-largeur au plus égale à $0{,}02$ (on se placera dans le cas le plus défavorable) ?
\end{enumerate}

\section*{Exercice 3 -- Lois théoriques et ajustement (5 points)}
\begin{enumerate}
  \item Le nombre $N$ de décès quotidiens dans un hôpital suit une loi de Poisson de paramètre $2$.
  Calculer $P(N = 0)$, $P(N \leqslant 2)$ et $P(N \geqslant 3)$.
  \item Les revenus annuels $X$ (en FCFA) des ménages d'une ville suivent une loi de Pareto de
  paramètres $x_0 = 50\,000$ et $\alpha = 2{,}5$ : $F(x) = 1 - \left(\dfrac{x_0}{x}\right)^{\alpha}$ pour $x \geqslant x_0$.
  \begin{enumerate}[label=\alph*)]
    \item Déterminer la densité de $X$.
    \item Calculer le revenu moyen (on justifiera son existence) et le revenu médian.
    \item On dispose d'un échantillon $(x_1,\dots,x_n)$, $x_0$ étant connu. Déterminer l'estimateur du
    maximum de vraisemblance de $\alpha$.
  \end{enumerate}
\end{enumerate}

\section*{Exercice 4 -- Série chronologique (5 points)}
Nombre trimestriel de naissances (en milliers) enregistrées dans une maternité :
\begin{center}
\begin{tabular}{lcccc}
\toprule
 & T1 & T2 & T3 & T4\\ \midrule
Année 1 & 10 & 14 & 18 & 12\\
Année 2 & 12 & 16 & 20 & 14\\
Année 3 & 14 & 18 & 22 & 16\\ \bottomrule
\end{tabular}
\end{center}
\begin{enumerate}
  \item Représenter la série et décrire ses composantes.
  \item Calculer les moyennes mobiles centrées d'ordre 4.
  \item En adoptant un modèle additif, calculer les coefficients saisonniers.
  \item Désaisonnaliser la série, puis prévoir le nombre de naissances du premier trimestre de l'année 4.
\end{enumerate}

\newpage
\begin{center}{\Large\bfseries Corrigé}\end{center}

\subsection*{Exercice 1}
\begin{enumerate}
  \item $\bar{x} = 3$, $\bar{y} = 4$ ; $\operatorname{Var}(x) = \frac{10}{5} = 2$ ; $\operatorname{Var}(y) = \frac{9{,}5}{5} = 1{,}9$ ;
  $\operatorname{Cov}(x,y) = \frac{9{,}7}{5} = 1{,}94$.
  \item $a = \frac{1{,}94}{2} = 0{,}97$, $b = 4 - 0{,}97\times3 = 1{,}09$ : $y = 0{,}97x + 1{,}09$.
  \item $r = \frac{9{,}7}{\sqrt{10\times9{,}5}} \approx 0{,}995$ : liaison linéaire positive très forte.
  \item $y(7) \approx 7{,}88$ millions ; extrapolation : suppose la tendance inchangée.
\end{enumerate}

\subsection*{Exercice 2}
\begin{enumerate}
  \item $\hat{p} = \frac{92}{400} = 0{,}23$.
  \item $0{,}23 \pm 1{,}96\sqrt{\frac{0{,}23\times0{,}77}{400}} = 0{,}23 \pm 0{,}041$, soit $[0{,}189\,;0{,}271]$.
  \item $z = \frac{0{,}23 - 0{,}20}{\sqrt{0{,}2\times0{,}8/400}} = \frac{0{,}03}{0{,}02} = 1{,}5 < 1{,}645$ :
  on ne rejette pas $H_0$ ; l'affirmation n'est pas établie au seuil de 5\,\%.
  \item $1{,}96\times\frac{0{,}5}{\sqrt n} \leqslant 0{,}02 \iff n \geqslant 2401$.
\end{enumerate}

\subsection*{Exercice 3}
\begin{enumerate}
  \item $P(N=0) = e^{-2} \approx 0{,}135$ ; $P(N\leqslant2) = e^{-2}(1 + 2 + 2) \approx 0{,}677$ ; $P(N\geqslant3) \approx 0{,}323$.
  \item a) $f(x) = \dfrac{\alpha\,x_0^{\alpha}}{x^{\alpha+1}}$ pour $x \geqslant x_0$, $0$ sinon.\\
  b) $E(X)$ existe car $\alpha > 1$ : $E(X) = \frac{\alpha x_0}{\alpha - 1} \approx 83\,333$ FCFA.
  Médiane : $\left(\frac{x_0}{M}\right)^{\alpha} = \frac12 \Rightarrow M = x_0\,2^{1/\alpha} \approx 65\,975$ FCFA.\\
  c) $\ln L(\alpha) = n\ln\alpha + n\alpha\ln x_0 - (\alpha+1)\sum\ln x_i$ ;
  $\frac{d\ln L}{d\alpha} = \frac n\alpha - \sum\ln\frac{x_i}{x_0} = 0$ ;
  $\hat\alpha = \dfrac{n}{\sum_{i=1}^n \ln(x_i/x_0)}$ (maximum car $\frac{d^2\ln L}{d\alpha^2} = -\frac{n}{\alpha^2} < 0$).
\end{enumerate}

\subsection*{Exercice 4}
\begin{enumerate}
  \item Tendance croissante (+2 par an et par trimestre) et saisonnalité de période 4 (pic au T3).
  \item Du T3 de l'année 1 au T2 de l'année 3 :
  $13{,}75$ ; $14{,}25$ ; $14{,}75$ ; $15{,}25$ ; $15{,}75$ ; $16{,}25$ ; $16{,}75$ ; $17{,}25$.
  \item Écarts $y - MM$ : T3 : $4{,}25$ ; T4 : $-2{,}25$ ; T1 : $-2{,}75$ ; T2 : $0{,}75$ (identiques chaque année),
  de somme nulle : $S_1 = -2{,}75$, $S_2 = 0{,}75$, $S_3 = 4{,}25$, $S_4 = -2{,}25$.
  \item Série CVS $= y - S_j$ : $12{,}75$ ; $13{,}25$ ; $13{,}75$ ; $14{,}25$ ; \dots (droite $13{,}75 + 0{,}5(t-3)$, $t=1$ pour le T1 de l'année 1).
  Pour $t = 13$ : tendance $18{,}75$, prévision $18{,}75 - 2{,}75 = 16$ milliers.
\end{enumerate}

\end{document}
